\subsection{Recovery Procedures}
\inlinedefinition[Recovery Manager] A recovery manager implements the recovery procedure, which depends on how the system functions.
It may or may not require both, one or none of UNDO and REDO.

\subsubsection{From the buffer cache's perspective}
\begin{itemize}
    \item \bi{Steal Policy}: An uncommitted transaction is allowed to overwrite the changes of a committed transaction in persistent storage.
          This logically means that we need to be able to undo transactions.
    \item \bi{Force Policy}: All changes made by a transaction must be in persistent storage before the transaction commits.
          This requires flushing all blocks with updates from the transaction and if not in-place, it will require being able to redo the transaction.
    \item \bi{Steal/No-Force}: UNDO/REDO, is the most common approach
\end{itemize}
If for updates tuples are locked, blocks may still be updatable.
Thus if the block is copied to storage, it is possible that some changes are committed while others are not.
If failures occur, there is no guarantee that the data in storage is 100\% consistent, thus requiring undo and redo.

\subsubsection{Recovery Procedure for UNDO/REDO}
\begin{itemize}
    \item READ: simply read the value from the block in the buffer cache
    \item WRITE: Create log entry (before and after image), append it to persistent log and write the after image to the block on the buffer cache.
    \item COMMIT: We write a persistent log entry indicating that the transaction has committed
    \item ABORT: For all updates, we restore the before image using the log entry
\end{itemize}

The Procedure:
\begin{enumerate}
    \item Start from the end of the log and work backwards
    \item Keep a list of UNDONE items and another for REDONE items
    \item Procedure terminates when all items are in either the UNDONE or REDONE list or we reach the beginning of the log.
    \item For each log entry, we look at the data item $x$ being accessed and if $x$ is in none of the two lists:
          \begin{itemize}
              \item If the log entry is of a committed transaction, apply the after image, add $x$ to the REDONE list
              \item If the log entry is of an aborted transaction, apply the before image, add $x$ to the UNDONE list
          \end{itemize}
\end{enumerate}
The procedure ignores the data stored on disk as it could correspond to uncommitted transactions, it only takes it as starting point for the recovery.
An alternative way of thinking about the procedure is as follows:
\begin{algorithm}
    \caption{Undo-Redo Recovery}
    \begin{algorithmic}[1]
        \Procedure{UndoRedoRecovery}{}
            \For{each item in the DB}
                \State Find last committed transaction modifying the item and REDO it
                \If{no committed transaction modified item}
                    \State Find the first aborted transaction that modified the item and \texttt{UNDO} the modification
                \EndIf
                \If{If no transaction has touched the item}
                    \State Value is correct (we assume that we start from a consistent state)
                \EndIf
            \EndFor
        \EndProcedure
    \end{algorithmic}
\end{algorithm}
This is not the way it's done in practice, because there typically are many more items than log entries.

\shade{ForestGreen}{Advantages} Only forced I/O are log records, buffer cache manager has a lot of freedom
(no flushing dirty pages, I/O on data is minimized and only triggered for replacement policies) and we allow writing dirty data to disk, simplifying buffer management.

However, recovery is more complicated and takes more time, but normal operations are only minimally affected.
Furthermore, queries are not affected since there is no forced I/O of data.
Transactions are affected, because we need to write every operation to the log.


\subsubsection{UNDO, no REDO}
\begin{itemize}
    \item READ: simply read the value from the block in the buffer cache
    \item WRITE: Create log entry (before image), append it to persistent log and write the after image to the block on the buffer cache.
    \item COMMIT: Flush all dirty values modified by transaction if still in cache. We write a persistent log entry indicating that the transaction has committed
    \item ABORT: For all updates, we restore the before image using the log entry
\end{itemize}
The Procedure:
\begin{enumerate}
    \item Start from the end of the log and work backwards
    \item Keep a list of UNDONE items
    \item Procedure terminates when all items are in either the UNDONE list or we reach the beginning of the log.
    \item For each log entry, we look at the data item $x$ being accessed and if $x$ is not in the UNDONE list and the transaction is aborted,
          we UNDO the changes by using the before image, adding $x$ to the UNDONE list
\end{enumerate}
It relies on the fact that all committed values are in persistent storage and therefore, they have not been lost.

However, it has forced I/O on all dirty blocks touched by the transaction when it commits,
log records no longer need to include after images, reducing log record size and the recovery procedure is shorter because it only involves undoing aborted transactions.

However, this trade-off does not pay off in practice, as we still need to write the log with every update and while it requires smaller log records,
it forces I/O on data blocks, which are often much larger.
And finally, flushing the buffer cache interferes with its operations.


\subsubsection{no UNDO, REDO}
\begin{itemize}
    \item READ: If the transaction did not write the item before, read the value from the block in the buffer cache. Otherwise, read the value from the temporary buffer
    \item WRITE: Create log entry (after image), append it to persistent log and write the after image to the block on the buffer cache.
    \item COMMIT: Apply all updates in the temporary buffer to the actual blocks. We write a persistent log entry indicating that the transaction has committed
    \item ABORT: Discard the temporary buffer
\end{itemize}
The Procedure:
\begin{enumerate}
    \item Start from the end of the log and work backwards
    \item Keep a list of REDONE items
    \item Procedure terminates when all items are in either the REDONE list or we reach the beginning of the log.
    \item For each log entry, we look at the data item $x$ being accessed and if $x$ is not in the REDONE list and the transaction is committed,
          we REDO the changes by using the after image, adding $x$ to the REDONE list
\end{enumerate}
It relies on the fact that there are never dirty blocks in the buffer cache. All data there is committed and is the last committed version.

It has forced I/O only on the log records and they are also smaller.
Furthermore, the recovery procedure is shorter, only involving redoing the last committed transaction that touched the item.
It is similar to snapshot isolation.


\subsubsection{No UNDO, no REDO}
This procedure is rarely used in conventional DBs. It does not require a log.
However, it requires to be able to write all changes made by a transaction to persistent storage in \bi{a single atomic action}.
It relies on directories to keep track of the changes.
\begin{itemize}
    \item READ: If the transaction did not write the item before, use current directory to find the latest committed copy.
          Otherwise, read the value from the shadow directory of that transaction to find the updated copy
    \item WRITE: Write to a buffer and add a pointer in the shadow directory for the transactions
    \item COMMIT: Create a full directory by merging the current one and the shadow directory of the transaction.
          Then, we swap the pointer indicating the latest committed directory.
    \item ABORT: Discard the buffer(s) and the shadow directory
\end{itemize}
There is no recovery procedure here, which is the exact point.

It is not used in practice, although some ideas can be applied partially.
One of the reasons for that is expensive pointer indirections for accesses and requires garbage collection and moves data all the time,
which poses a challenge for clustered indexes or hash clustered tables.
